\section{Ideas 8--12: how the finite input becomes a proof}\label{sec:computation}
\subsection{Idea 8: search for weights, not just large examples}
A useful search describes a small family by a vector of counts:
its size, available swaps, endpoint types and selected smaller slices.
The desired inequality is linear in unknown weights on those counts.
A linear program can seek weights making all current constraints
valid. Tight or failing families suggest which additional statistics
or corrections may be useful. This explains the purpose of the
five-type potential and the exceptional correction.

The discovery record reports floating-point linear programming and
an integer correction. It does not document every possible intermediate
heuristic. Trying candidates, keeping useful choices, and examining
counterexamples is a legitimate discovery strategy, but its selection
history is not a theorem premise. The final proof freezes explicit
rational data and proves every required assertion about those data.
The coefficient \(5/8\), rather than a small-domain maximum size,
is what feeds the asymptotic bound.
\outside{The historical search method and this explanatory description
of a possible iterative workflow are not Lean theorems. Lean proves
properties of the fixed data that the search eventually supplied.}

\subsection{Idea 9: exhaustive constraint generation}
An extreme-law-plus-pair family is specified by one of six extreme
laws on each triple and one of three choices on each pair: absent,
first orientation, second orientation. Regard every choice as the
subset of permutations it permits. The full solution set is the
intersection of the selected subsets. Therefore exploring all such
choices covers every family under consideration, including
non-maximal ones. Conflicting choices may give the empty family.
Maximal-domain classifications alone would not prove this coverage.

Here is the basic checker in mathematical pseudocode:
\[
\begin{split}
 \operatorname{Check}(S,[])&=\operatorname{Leaf}(S),\\
 \operatorname{Check}(S,[\mathcal O_1,\ldots,\mathcal O_m])
  &=\bigwedge_{O\in\mathcal O_1}
       \operatorname{Check}(S\cap O,[\mathcal O_2,\ldots,\mathcal O_m]).
\end{split}
\]
Empty intersections can be accepted immediately when the target
inequality is zero on the empty family. Correctness is induction on
the remaining groups: each actual choice selects one branch, and its
final solution set is the intersection at that leaf. Thus a true
checker result proves the property for every choice. Deduplicating
identical states or relabelling to a proved representative is safe;
sampling only promising branches would not be.

\lean{\path{CDCert.IsChoice}, \path{domainOf}, \path{allows},
\path{CDCert.ListOrder.groupsEP} describe the constraints. The
semantic bridge \path{CDCert.Bridge.isChoice_groupsEP_iff}
shows exactly which law/pair choices the generator represents.
The soundness files include \path{CDCert/K6/Sound.lean},
\path{DfsPSem.lean} and \path{FamSem.lean}; the specialised depth-first
search checks are connected to the abstract choices by proved lemmas,
not by accepting the output of a native program.}

\subsection{Coverage on five, six and seven alternatives}
At five alternatives the fixed table covers 1160 nonempty types.
For any defining system on five labels, restrict its constraints to
four labels, classify that full four-label solution set, then add back
the constraints involving the new label. Restricting \emph{constraints}
is essential: projecting the permitted orders can lose the full
constraint description. The complete proof is
\path{CDCert.K6.five_cover}; it gives a relabelling to one of the
1160 representatives. The initial four-level classification is also
proved in the project.

For a six-family, restrict its defining constraints to five labels.
That full five-family appears in the table, even if it is larger than
the projection of the six-family. Reintroduce every possible new
extreme triple law and every new pair choice. One branch is the
original six-family. Kernel-evaluated checks certify either the
regular inequality or membership in the complete exceptional list.
The theorem \path{CDCert.K6.six_sem} is the universal result; its
semantics identify the regular residual with
\[
 ds(H)-\eta d\bigl(|H|+4r_b(H)+4r_t(H)\bigr).
\]
Thus a nonnegative residual gives F2 after substituting the correction
\(p\), while the stored exceptional types provide F1 and the remaining
seven-level cases. All table fields are proved to be the sizes,
adjacencies, slice values and type indicators they claim to encode.

For seven alternatives, F3 requires every family with an exceptional
six-slice other than \(R_8\). Relabel that slice so its removed
endpoint is label 6 and its remaining labels match the exact stored
six-mask. Enumerate every defining constraint choice compatible with
that exact endpoint slice. This is done for all thirteen relevant
six-types at both top and bottom, giving 26 oriented seed families.
The theorems \path{seven_fam_0}, \ldots, \path{seven_fam_25} prove
\(\Lambda\geq\eta|J|\) under the corresponding exact-slice hypothesis.
\path{relabel_slice}, \path{seven_bot} and \path{seven_top} transport
these results to every relevant slice at any label. Relabelling
preserves cardinality and \(\Lambda\), so it loses no case.

The formal finite enumeration has 16888 final states. These are
counts of states in this proof, not a value of \(f(7)\). The independent
native implementation uses different reductions and has different
counts. Agreement of those literal counts is not needed: completeness
and the same universal inequality are what matter.

\subsection{Idea 10: exceptions are mathematical information}
The regular estimate intentionally permits a small specified list
of failures. Thirteen types are covered by the oriented seven-family
checks. The remaining \(R_8\) case is proved by a hand argument:
its common first and last labels force a receiving slice with a
specific extra five-slice. Different deficits have distinct
receiver/label pairs. Section~\ref{sec:local} proves this charging
argument in full. The computer-supported proof is therefore a
combination of exhaustive finite checks and a general structural
explanation of the one remaining exception, not merely a list of
successful inequalities.

\subsection{Idea 11: exact certificates and a small trusted kernel}
All rational coefficients have the common denominator \(d\).
The finite checks use integer residuals and prove nonnegativity
exactly. In Lean, the large generated checks use kernel reduction,
including \texttt{decide +kernel}; they do not use
\texttt{native\_decide}, a floating-point tolerance or an axiom saying
that an external checker succeeded. Generic soundness theorems turn
those finite evaluations into the universal mathematical statements.
The full source project includes every generated branch and literal
table used in that chain. It requires no unprovided generator to
understand the statements or to recompile their proofs.

\subsection{Idea 12: prove that the representation preserves the mathematics}\label{sec:bridge}
The semantic definitions use Boolean comparison relations; the finite
checker uses permutations as lists. This difference must be proved
harmless. Sorting a semantic order gives a list \(\operatorname{toList}(r)\)
containing every label exactly once. Reading precedence in that list
recovers exactly the original comparison relation. The inverse
\(\operatorname{ofList}\) constructs a strict total order from a
permutation. The two conversions are inverse, so no order is lost
and no duplicate is introduced. A never-law or fixed pair holds for
an order exactly when the corresponding list option accepts its list.
Consequently full solution sets, their cardinalities and adjacent
swaps agree on both sides. Endpoint deletion also uses the explicit
canonical relabelling already described.

\lean{\path{CDCert.Bridge.before_toList}, \path{ofList_toList},
\path{toList_ofList}, \path{toList_injective},
\path{lawOpt_toList}, \path{pairOpt_toList},
\path{mem_extremeFullDomain_iff}, \path{InF.toList_domain},
\path{toList_swapLetter} and \path{card_adj_of_choice} establish
these connections. The final theorem \path{CDCert.Bridge.InF.density}
uses those equalities, \path{CDCert.K6.adj_le}, the semantic fibre
identity and gap zero. Its conclusion is about actual semantic
orders, not merely a model that resembles them.}

The native Python/C++ programs are useful independent replay tools,
but their implementation correctness, timings, literal counts and
search history are not propositions proved by this Lean project.
The main theorem relies on the separately proved Lean enumeration
and its semantic bridge.
